\begin{proof}
Write $\mathcal J_{\mathrm{eucl}}$ and
$\mathcal G_{\mathrm{eucl},0}$ for the exact comparison and generator
sheaves of
\Cref{thm:p071-infinite-star-euclidean-collapse} with one Euclidean
coordinate.

\medskip\noindent
\emph{Locality and decreasing closure.}
All displayed comparison conditions are local.  The contraction argument
in the proof of the cited theorem shows that the junction inequality is
unchanged when a full branch becomes exceptional after a restriction; the
same argument works at $s=0$ because the homothety in the $s$-coordinate
preserves the half-interval.  Joint convexity glues on branch products, the
full-endpoint trace glues on a terminal branch collar, and the
boundary-minimum condition glues on the connected endpoint component.
Thus $\mathcal J_+$ is a sheaf.

It is closed under pointwise decreasing sequential limits.  Joint
extended-real convexity, the boundary-minimum inequality, and the junction
inequality pass to the limit exactly as in the Euclidean product theorem.
For the latter, if the limiting central value is finite, all finitely many
values in the displayed tuple are finite, so ordinary termwise passage is
legitimate.  The full-endpoint trace passes by the monotonicity argument in
that theorem.  At a node, compactness of a closed endpoint interval and the
cofinite-branch topology give one countable activity set for each row; the
union of those sets over the sequence is still countable and works for the
limit.

Every generator belongs to $\mathcal J_+$.  On the finite stratum this is
the generator inclusion of the Euclidean product theorem.  On its initial
constant collar, the boundary minimum is immediate.  Beyond that collar,
joint convexity makes the generator nondecreasing in $s$, because its
$s$-derivative starts at zero and is nondecreasing along every vertical
line.  The same derivative-balance calculation as in the cited theorem
gives \eqref{eq:p071-star-finite-collar-jensen}, now allowing some or all
of the $s$-coordinates to be zero.  Minimality therefore gives
\begin{equation}
                         \mathcal P_+\subseteq\mathcal J_+.
\label{eq:p071-star-finite-collar-minimality}
\end{equation}

\medskip\noindent
\emph{A monotone-coordinate refinement of the Euclidean theorem.}
We record the exact refinement needed below.  If
$v\in\mathcal J_{\mathrm{eucl}}$ is nondecreasing in its Euclidean
coordinate, then the fixed-neighbourhood approximating sequence supplied
by \Cref{thm:p071-infinite-star-euclidean-collapse} may be chosen as
\begin{equation}
             g_{n+1}\leq g_n,\qquad g_n\downarrow v,
             \qquad \partial_s g_n\geq0.
\label{eq:p071-star-finite-collar-monotone-refinement}
\end{equation}

Here are the details, including the extended-real star endpoints.  In the
partial-Legendre step of the proof of the Euclidean theorem, work on the
same compact inner interval in the $s$-variable.  A finite convex function
which is nondecreasing in $s$ is the supremum, on every smaller interval,
of its supporting affine functions whose $s$-slopes are nonnegative.
Accordingly, every finite dual sample in that proof may be chosen with
nonnegative $s$-component.  The scalar exact-star truncations and scalar
star smoothing do not alter that component.  Each lifted smooth function
therefore has nonnegative $s$-derivative, and a smooth convex
coordinatewise nondecreasing regularized maximum preserves this property.

At a full star endpoint, the proof first replaces $v$ by
\[
                 H_m(x,s)=\max\{v(x,s),v(x_1,s)-m\}
\]
on a compact terminal core, for one fixed finite branch point $x_1$.
Both entries are nondecreasing in $s$, so $H_m$ is also nondecreasing.
The same nonnegative-slope dual sampling and endpoint smoothing therefore
apply.  If a node central value is $-\infty$, the cited proof uses constant
approximants after shrinking, which already have zero $s$-derivative.  At
an ordinary branch point, extend a finite jointly convex nondecreasing
function to a slightly larger convex set as a supremum of supporting
affine functions with nonnegative $s$-slope, take a finite maximum on the
chosen compact core, and mollify; this gives the same property.  Finally,
the summable tail constants used to impose exact decrease do not change any
derivative.  These observations prove
\eqref{eq:p071-star-finite-collar-monotone-refinement} in every stratum.

\medskip\noindent
\emph{Convex extension across the finite collar endpoint.}
Let $u\in\mathcal J_+$ near a point $(x_0,0)$.  Shrink to a displayed
product neighbourhood $V_X\times[0,r)$.  For each fixed $x$, joint
convexity and
\eqref{eq:p071-star-finite-collar-boundary-minimum} show that
$s\mapsto u(x,s)$ is nondecreasing.  Define on
$V_X\times(-r,r)$
\begin{equation}
             \widetilde u(x,s)=u\bigl(x,s^+\bigr),
             \qquad s^+=\max\{s,0\}.
\label{eq:p071-star-finite-collar-extension}
\end{equation}
We claim that $\widetilde u\in\mathcal J_{\mathrm{eucl}}$ after one
harmless shrinking.

On each ordinary branch product, this follows from the standard convex
composition rule: the inner map fixes the branch coordinate and replaces
$s$ by the convex function $s^+$, while the outer jointly convex function
is nondecreasing in its second variable.  The exact full-branch endpoint
trace and the common countable activity condition pass unchanged through
\eqref{eq:p071-star-finite-collar-extension}.

It remains to check the product junction inequality.  Take a junction
tuple for $\widetilde u$ with central coordinate $s_0$ and branch
coordinates $s_i$, using reciprocal-distance weights whose sum is
$\Lambda$.  Put
\[
             q=\Lambda^{-1}\sum_i\lambda_i s_i^+,
             \qquad \lambda_i\in
               \{1/r_b:b\in F\}\cup\{1/\rho_a:a\in A\}.
\]
The original barycentre relation gives
\begin{equation}
                         q\geq s_0^+.
\label{eq:p071-star-finite-collar-positive-part-barycentre}
\end{equation}
After shrinking the ambient interval once, $q$ remains in $[0,r)$.
Since $c$ is nondecreasing, the junction inequality for $u$, applied at
the centre coordinate $q$ and the sample coordinates $s_i^+$, gives
\[
 \Lambda\widetilde c(s_0)
   =\Lambda c(s_0^+)
   \leq\Lambda c(q)
   \leq\sum_i\lambda_i u_i(s_i^+)
   =\sum_i\lambda_i\widetilde u_i(s_i).
\]
This is precisely the required junction inequality for the extension.
The claim follows.

Apply the monotone-coordinate refinement to obtain, on one fixed smaller
product neighbourhood $V_X'\times(-\delta,r')$,
\begin{equation}
        g_{n+1}\leq g_n,qquad g_n\downarrow\widetilde u,
        \qquad \partial_sg_n\geq0,
\label{eq:p071-star-finite-collar-euclidean-approximants}
\end{equation}
with $g_n\in\mathcal G_{\mathrm{eucl},0}$.

\medskip\noindent
\emph{Restoring exact constant collars.}
Choose a smooth convex nondecreasing function
$\phi:[0,+\infty)\to\mathbf R$ which is constant near zero, agrees with
the identity for $q\geq2$, and satisfies $\phi(q)\geq q$.  Such a function
is obtained by choosing a smooth nondecreasing derivative which is zero
near zero and one on $[2,+\infty)$, and then choosing the integration
constant so that $\phi(2)=2$.  For $a>0$, put
\begin{equation}
                         h_a(s)=a\phi(s/a).
\label{eq:p071-star-finite-collar-clamp}
\end{equation}
Then $h_a$ is smooth, convex, nondecreasing, constant on an initial
interval, and equal to $s$ for $s\geq2a$.  Moreover
\begin{equation}
              0<a'<a\quad\Longrightarrow\quad h_{a'}\leq h_a.
\label{eq:p071-star-finite-collar-clamp-order}
\end{equation}
Indeed, with $q=s/a$,
\[
       \frac{\partial}{\partial a}\bigl(a\phi(s/a)\bigr)
          =\phi(q)-q\phi'(q)\geq0;
\]
the last function is nonincreasing in $q$ and equals zero for $q\geq2$.

After shrinking $V_X'$ once more, choose a compact core $D$ contained in
the fixed domain of the $g_n$ and containing the closure of $V_X'$ in its
interior.  At a node, take a compact closed standard star core; at an
ordinary branch point, take a compact branch interval; at a full branch
endpoint, take a compact terminal branch interval including that endpoint.
Fix $0<r''<r'$.  Recursively
choose $a_n\downarrow0$ so rapidly that $a_{n+1}<a_n/2$, the image of
$[0,r'')$ under $h_{a_n}$ lies in $(-\delta,r')$, and
\begin{equation}
 \sup_{(x,s)\in D\times[0,r'']}
    \bigl|g_n(x,h_{a_n}(s))-g_n(x,s)\bigr|<\frac1n.
\label{eq:p071-star-finite-collar-uniform-clamp-error}
\end{equation}
This is possible because each real-valued generator $g_n$ is uniformly
continuous on a slightly larger compact core and $h_a\to\operatorname{id}$
uniformly on compact half-intervals.

Set
\begin{equation}
                     G_n(x,s)=g_n(x,h_{a_n}(s)).
\label{eq:p071-star-finite-collar-final-approximants}
\end{equation}
The function $G_n$ is an exact product generator.  Joint convexity follows
from convex monotone composition.  It is independent of $s$ on an initial
collar.  At the star node, its finite activity set is that of $g_n$ and
\[
 \sum_b\partial_tG_{n,b}(0,s)
   =\sum_b\partial_tg_{n,b}(0,h_{a_n}(s))=0.
\]
The terminal-collar property at a full star endpoint is also unchanged.

Equations
\eqref{eq:p071-star-finite-collar-monotone-refinement} and
\eqref{eq:p071-star-finite-collar-clamp-order} give
\[
 G_{n+1}(x,s)
 =g_{n+1}(x,h_{a_{n+1}}(s))
 \leq g_n(x,h_{a_{n+1}}(s))
 \leq g_n(x,h_{a_n}(s))=G_n(x,s).
\]
By \eqref{eq:p071-star-finite-collar-uniform-clamp-error} and
\eqref{eq:p071-star-finite-collar-euclidean-approximants},
$G_n(x,s)\to u(x,s)$ pointwise on
$V_X'\times[0,r'')$, including every extended-real value on a full star
endpoint.  Thus every endpoint germ in $\mathcal J_+$ belongs to
$\mathcal P_{+,1}$ on one fixed smaller product neighbourhood.

At a point with positive $S_+$-coordinate, the assertion is exactly
\Cref{thm:p071-infinite-star-euclidean-collapse}.  Hence
$\mathcal J_+\subseteq\mathcal P_{+,1}$ everywhere.  Together with
\eqref{eq:p071-star-finite-collar-minimality} and the tautological
inclusion $\mathcal P_{+,1}\subseteq\mathcal P_+$, this proves
\eqref{eq:p071-star-finite-collar-collapse} and the fixed-neighbourhood
assertion.
\end{proof}
